\documentclass{book}
\usepackage[backend=biber,natbib=true,style=authoryear]{biblatex}
\addbibresource{/home/nguyen/1_NQBH/math/bib.bib}
\usepackage[utf8]{vietnam}
\usepackage{float}
\usepackage{graphicx}
\usepackage[colorlinks=true,linkcolor=blue,urlcolor=red,citecolor=magenta]{hyperref}
\usepackage{amsmath,amssymb,amsthm}
\allowdisplaybreaks
\numberwithin{equation}{section}
\newtheorem{assumption}{Assumption}[section]
\newtheorem{lemma}{Lemma}[section]
\newtheorem{corollary}{Corollary}[section]
\newtheorem{definition}{Definition}[section]
\newtheorem{proposition}{Proposition}[section]
\newtheorem{theorem}{Theorem}[section]
\newtheorem{notation}{Notation}[section]
\newtheorem{remark}{Remark}[section]
\newtheorem{example}{Example}[section]
\newtheorem{ques}{Question}[section]
\newtheorem{problem}{Problem}[section]
\newtheorem{conjecture}{Conjecture}[section]
\usepackage[left=0.5in,right=0.5in,top=1.5cm,bottom=1.5cm]{geometry}
\usepackage{fancyhdr}
\pagestyle{fancy}
\fancyhf{}
\addtolength{\headheight}{0pt}% obsolete
\lhead{\itshape \scriptsize \chaptername~\thechapter}
\rhead{\itshape \scriptsize \nouppercase{\leftmark}} %\nouppercase !
\renewcommand{\chaptermark}[1]{\markboth{#1}{}}
\cfoot{\thepage}

\title{\href{http://tramdoc.vn/}{Trạm Đọc}}
\author{Collector: Nguyễn Quản Bá Hồng}
\date{\today}

\begin{document}
\maketitle
\setcounter{secnumdepth}{6}
\setcounter{tocdepth}{6}
\tableofcontents

%------------------------------------------------------------------------------%

\chapter{\href{http://tramdoc.vn/tin-tuc/phuong-trinh-hanh-phuc-cua-einstein-nzGgvW.html}{Phương Trình Hạnh Phúc của Einstein}}

\textbf{Chuyện gì cũng có công thức khoa học.}

%
Sống trên đời làm gì cũng cần lý do.

Từ lâu loài người đã đặt ra câu hỏi: chúng ta tồn tại để làm gì?

Một vài người đàn ông thích theo đuổi phụ nữ, một vài người phụ nữ thích theo đuổi tiền tài, một số khác thì theo đuổi những cuộc vui chuếnh choáng, nhưng tựu chung lại tất cả mọi người đều mưu cầu cùng một thứ: hạnh phúc đích thực.
\begin{quotation}
    ``\textit{Hạnh phúc là ý nghĩa và mục đích của cuộc đời, là đích đến và kết cuộc cho sự tồn tại của nhân loại}.'' - Aristotle
\end{quotation}
Vì vậy để đạt được hạnh phúc ta thử tham khảo xem những nhà thông thái đã làm gì.

Einstein, người được xem là nhà khoa học vĩ đại nhất lịch sử loài người, đã viết ra một phương trình hạnh phúc:

\section{Phương trình hạnh phúc của Einstein}
Trong những ngày cuối đời, Dukas - trợ lý của Einstein - miêu tả rằng ông nằm trên giường bệnh \textit{trong đau đớn và không thể ngẩng đầu lên}.

%
Tuy nhiên trong ngày tiếp theo, chỉ cách thời khắc sinh tử không đầy 24h, Einstein nhờ Dukas \textit{lấy giúp ông mắt kính, giấy và bút chì để ông ghi vội một vài công thức}.

%
\textit{Ông ấy viết rất lâu}, Sử gia Walter Isaacson cho biết, \textit{cho đến khi cơn đau trở nặng không thể chịu đựng được nữa, ông ấy ngủ thiếp đi} - một giấc ngủ vĩnh hằng.

Einstein qua đời trong khi vẫn làm điều mà ông yêu thích nhất - làm việc.

Thật không ngoa khi nói hoàn cảnh làm bật lên nhân cách của một con người.

%
Einstein từng nói: \textit{Thiên tài chỉ có 1\% là bẩm sinh và 99\% là nỗ lực}.

Trong suốt chiều dài lịch sử không có ai trở nên vĩ đại hoàn toàn dựa vào vận may.

Ông cho rằng: \textit{Chỉ có những người chăm chú vào duy nhất một vấn đề mới có thể gặt hái được thứ mà ta thường gọi là `thành quả'}.

%
Chỉ cho tôi một người vĩ đại và tôi sẽ chỉ cho bạn một kẻ điên cuồng.

Có gì gần với ``sự vĩ đại'' hơn sự ám ảnh điên cuồng kia chứ?

%
Vì những lý do trên, có người đến hỏi Einstein bí quyết để có một cuộc sống hạnh phúc là gì.

Người phỏng vấn nghĩ rằng ông sẽ đáp lại rất dài dòng, nhưng ông chỉ trả lời ngắn gọn:
\begin{quotation}
    ``\textit{Nếu bạn muốn sống một cuộc đời hạnh phúc, hãy gắn nó với một mục đích, đừng gắn nó với bất kỳ ai hay bất cứ điều gì.}''
\end{quotation}
Và đó là phương trình hạnh phúc của Einstein.

\section{Công việc - Ngôi nhà an toàn trước mọi khó khăn của cuộc đời}
Những bộ óc vĩ đại thường suy nghĩ giống nhau.

Chẳng lạ gì khi cả Newton và Einstein - 2 bộ óc vĩ đại của loài người - khi đối mặt với khó khăn trong cuộc sống đều dùng chung một cách thức để vượt qua.

%
Năm 1665, dịch hạch hoành hành khắp London.

Lịch sử loài người chưa từng chứng kiến cơn bùng phát dịch nào khủng khiếp đến thế.

Sức tàn phá của nó không thua gì Covid-19 hiện nay, khiến các trường học và cửa hàng phải đóng cửa.

%
Khi thế giới đang chìm trong hoảng loạn, Isaac Newton đã dùng đến công thức hạnh phúc của mình.

Dùng như thế nào?

Đơn giản là ông gắn cuộc đời mình vào một mục tiêu trừu tượng chứ không phải một con người hay bất kỳ vật thể hữu hình nào khác.

Trong vòng 18 tháng sống cách ly tại một trang trại, ông đã khởi đầu một số công trình nghiên cứu mà theo đánh giá của các nhà chuyên môn là ``làm thay đổi thế giới''.

%
Sau này, khi được hỏi làm cách nào mà ông tìm ra luật vạn vật hấp dẫn, Newton đã trả lời:
\begin{quotation}
    \textit{Tôi nghĩ về nó mọi lúc mọi nơi}.
\end{quotation}
Tóm lại, vì Newton gắn liền cuộc đời mình vào một mục tiêu lý tưởng - một điều không bao giờ thay đổi, chứ không phải cuộc đời, vốn luôn rất khó đoán, nên đã bảo vệ bản thân trước mọi biến động của cuộc sống.

%
Khi vợ của Einstein, bà Elsa, qua đời - ông rơi vào khủng hoảng.

Theo sử gia Isaacson, bà không chỉ là vợ mà còn giống như một người mẹ đối với Einstein.

%
\textit{Bà nói cho ông biết khi nào thì nên ăn và nơi nào thì nên đi}, Isaacson miêu tả, \textit{Bà soạn hành lý cho ông rồi chuẩn bị sẵn tiền tiêu vặt}.

\textit{Ở những nơi công cộng, bà luôn che chở cho ông và gọi ông là `Giáo sư'}.

%
May thay cho Einstein, phương trình hạnh phúc đã giúp ông vượt qua sự mất mát ấy.

\textit{Khi nào tôi còn có thể tiếp tục làm việc, tôi không nên và sẽ không bao giờ than phiền, bởi vì công việc là nguồn sống duy nhất của đời tôi}.

\section{Phương trình Hạnh phúc $=$ Dùng tâm trí để mơ $+$ Dùng tất cả sức lực để theo đuổi giấc mơ đó}
Khi một nhà sư được hỏi ông có cảm thấy ổn không khi từ bỏ những thú vui trần thế, ông trả lời rằng: \textit{Tôi không cảm thấy tốt hơn, nhưng cũng không còn cảm thấy tệ nữa}.

%
Công thức để có một cuộc sống hạnh phúc mà Einstein đúc kết lại là:
\begin{quotation}
    ``\textit{Hạnh phúc giống như một siêu mẫu chuẩn bị kỹ càng để  bước lên thảm đỏ. Tuy nhiên sự thỏa mãn lại là lúc khi cô người mẫu ấy trở về nhà và gỡ bỏ lớp hóa trang}.''
\end{quotation}
Nói một cách khác, giống như Socrates từng nhận xét, \textit{Người giàu có nhất là người hài lòng với ít thứ nhất}.

%
Những điều tốt đẹp nhất trong cuộc sống không thể tốn tiền để đạt được cũng  không phải là vật chất hữu hình.

Liệu có ai từng chạm vào tình yêu?

Liệu có ai từng nhìn thấy hòa bình?

Liệu có ai từng ngửi thấy mơ ước?

Dưới đây là yếu tố cốt lõi trong phương trình Hạnh phúc của Einstein.

%
Trong tất cả các loài động vật, con người là loài duy nhất có ước mơ.

Đó là một món quà diệu kỳ được gọi là ``tâm thức''.

Chúng ta có thể tưởng tượng ra một viễn cảnh hoàn hảo rồi dồn hết toàn bộ năng lượng từ trái tim và linh hồn để biến điều đó thành sự thật.

%
Thông điệp ở trên đã được một nhân vật trong vở kịch kinh điển \textit{The Secret of Freedom} (tạm dịch: Bí mật của Tự do) tóm tắt như sau:

%
\textit{Điều làm nên một con người chính là tâm thức của anh ta}.

\textit{Những thứ khác thì ngựa hay heo cũng có}.

%
Tuy nhiên, người ta thường nói những thứ cho không ít được ai quý trọng.

%
Từ lúc ra đời, mỗi con người đều được ban tặng nhiều phẩm hạnh tốt đẹp, ví dụ như tình yêu và lòng tin.

Có lẽ vì thế mà chúng ta đã đánh giá thấp khả năng mơ đến một viễn cảnh xứng đáng chứ đừng nói đến khả năng biến ước mơ đó thành mục đích của cuộc đời.

Ước mơ, giống như tình yêu và lòng tin, là những phước lành ta được ban tặng miễn phí từ khi mới lọt lòng.

Có lẽ vì vậy mà ta đã xem nhẹ khả năng mơ ước đến những mục tiêu đáng giá, chứ đừng nói đến khả năng kết nối cuộc đời với những mục tiêu như vậy.

%
Tôi xin được lấy ví dụ một người bạn của mình và cách cô ấy dùng phương trình hạnh phúc của Einstein để vượt qua giông bão.

%
Một người bạn rất thân thời đại học của tôi đã trải qua một nỗi đau khôn tả.

Hôn phu của cô ấy, người cô quen từ thời trung học, phản bội cô và đi theo một người phụ nữ khác.

Vậy mà bằng cách nào đó cô ấy đã vượt qua được cú sốc đó, giống hệt người phụ nữ trong bài thơ \textit{Phenomenal Woman} của Maya Angelou.

%
Cô đã làm điều đó bằng cách nào?

%
Cô tự vực bản thân dậy khỏi nghi ngờ và sợ hãi, làm sạch trái tim  mình, rồi dồn hết thời gian và năng lượng vào việc học.

Khoảng một năm sau cô được nhận vào Trường Luật Loyola.

%
Ngày hôm nay, không những có một cuộc hôn nhân viên mãn mà cô còn là một công tố viên thành đạt.

Bạn của tôi đã áp dụng phương trình Hạnh phúc của Einstein vào cuộc sống của cô ấy.

\textsf{Hạnh phúc cũng cần có phương trình của mình.}

\section{Kết luận}
Vì Einstein là một nhà toán học, ông xem những bí quyết dẫn đến thành công là thành tố trong một phương trình.

\textit{Bạn phải nắm quy tắc của trò chơi}, ông kết luận, \textit{và sau đó bạn phải chơi giỏi hơn những người khác}.

%
Dường như quy tắc của Trò Đời là chúng ta - sinh vật duy nhất có khả năng tưởng tượng - phải sử dụng nó để hướng tới một mục đích trừu tượng, hay còn được gọi là ``ước mơ.''

%
Có thể những lý do trên lý giải được vì sao Einstein cho rằng \textit{trí tưởng tượng quan trọng hơn kiến thức}.

Cho dù gắn liền cuộc sống với giấc mơ không hẳn là lựa chọn huy hoàng nhất, nhưng ông cho rằng đó là con đường dễ đi nhất trong cuộc đời này.

%
Tất cả chúng ta đều không biết vào giây phút kế tiếp chuyện gì sẽ xảy ra - đó có thể là vinh quang, cũng có thể là bi kịch.

Một cách sống khôn ngoan là xem cuộc đời như một chiếc thuyền và biến ước mơ thành mỏ neo có thể giúp ta đứng vững qua giông bão.

%
Tóm lại, phương trình Hạnh phúc của cuộc đời chính là:
\begin{quotation}
    ``\textit{Tìm ra thứ bạn làm tốt nhất; sau đó tìm cách để tập trung tối đa vào nó; sau đó tìm người trả tiền cho bạn vì bạn đã dành toàn bộ công sức để hoàn thiện kỹ năng đó}.''
\end{quotation}

\begin{flushright}
    Vy Vũ/Theo \href{https://medium.com/mind-cafe/einsteins-formula-for-a-happy-life-b29aff61a9c7}{Medium}
\end{flushright}

%------------------------------------------------------------------------------%

\chapter{\href{http://tramdoc.vn/tin-tuc/9-thoi-quen-ky-la-tao-nen-phong-cach-cua-cac-nha-van-vi-dai-nmwejW.html}{9 Thói Quen Kỳ Lạ Tạo Nên Phong Cách của Các Nhà Văn Vĩ Đại}}

\textbf{Bên cạnh tài năng thiên bẩm, những con người ấy còn sở hữu lòng nhiệt thành và đam mê với nghề. Các thói quen có phần kỳ quặc của họ sẽ chắp cánh cho ngôn từ bay bổng trên trang giấy. Có thể, chỉ một cái nhìn thoáng qua về cuộc sống thường ngày cũng có thể giúp họ tìm thấy cảm hứng sáng tạo.}

%
Trong cuộc chiến với trang giấy trắng, tác giả cần phải có một chiến lược vững chắc.

Ý niệm này không chỉ áp dụng cho những nhà văn trẻ mà còn với cả \href{http://tramdoc.vn/tin-tuc/10-tac-pham-kinh-dien-tu-giai-nobel-van-hoc-da-duoc-dich-sang-tieng-viet-nyalW.html}{những tượng đài của văn học}.

Trên con đường tạo ra một kiệt tác, những bậc thầy ngôn từ cũng phải chật vật với sự kỳ vọng vào việc tạo ra tác phẩm tốt nhất và khó khăn trong việc tìm kiếm cảm hứng cho bản thân mình.

%
Bên cạnh tài năng thiên bẩm, những con người ấy còn sở hữu lòng nhiệt thành và đam mê với nghề. Các thói quen có phần kỳ quặc của họ sẽ chắp cánh cho ngôn từ bay bổng trên trang giấy.

Có thể, chỉ một cái nhìn thoáng qua về cuộc sống thường ngày cũng có thể giúp họ tìm thấy cảm hứng sáng tạo.

\section{Nằm để viết: Tư thế lý tưởng mang đến nguồn cảm hứng bất tận}
Với một số nhà văn, tư thế này giúp họ giải phóng năng lượng và có thể tập trung hơn.

Sự êm ái của chiếc giường sẽ đem lại cảm hứng và những từ ngữ hoàn hảo.

Có thể kể đến một số nhà văn sở hữu thói quen này như Mark Twain, George Orwell, Edith Wharton, Woody Allen và Marcel Proust.

Giường và ghế sofa là nơi đã giúp họ tạo ra vô số trang tuyệt phẩm.

Nhà văn, nhà viết kịch người Mỹ Truman Capote tự gọi mình là ``một nhà văn nằm'', bởi vì ông không thể suy nghĩ và viết ở những tư thế khác.

%
\textsf{Mark Twain trên chiếc giường của mình.}

\section{Đứng để sáng tác}
Ngược lại với thói quen trên, đứng viết dường như không còn quá xa lạ với những nhà văn theo đuổi dòng sách có yếu tố kịch tính, giật gân.

Danh sách này bao gồm \href{http://tramdoc.vn/tin-tuc/25-cuon-sach-ma-nguoi-dan-ong-nao-cung-nen-doc-it-nhat-mot-lan-trong-doi-nZN4aW.html}{Ernest Hemingway}, Charles Dickens, Virginia Woolf, Lewis Carroll và Philip Roth.

Nguồn cảm hứng cho những tác phẩm lạ, có phần đặc biệt này đến từ chiếc bàn đứng.

Phương pháp này cũng đặc biệt có lợi cho những người quan tâm tới sức khỏe, bởi vì khoa học đã chứng minh đứng viết sẽ có lợi cho cột sống của chúng ta hơn.

\section{Những tờ giấy nhớ}
Vladimir Nabokov, tác giả của cuốn Lolita, Pale Fire và Ada, đã có một cách tiếp cận hết sức đặc biệt với việc sáng tác.

Ông đã tạo ra những điều tuyệt vời chỉ với vài tờ giấy nhớ nhỏ nằm trong chiếc hộp mỏng.

Phong cách khác thường này giúp tác giả có thể tạo dựng nên hàng loạt bối cảnh khác nhau của cuốn tiểu thuyết, và sau đó có thể thay đổi vị trí theo ý muốn.

Và  một điều đặc biệt khác là, Nabokov luôn có những tờ giấy nhớ dưới gối để phòng khi có một ý tưởng bất chợt nảy ra trong đầu, ông sẽ lưu lại ngay lập tức. 

Biết đâu phương pháp của Nabokov cũng có ích với bạn thì sao?

\section{Sử dụng thủ thuật màu sắc}
Tác giả của những cuốn tiểu thuyết như  ``The Three Musketeers'' - Ba chàng lính ngự lâm và ``Count of Monte Cristo'' - Bá tước Monte Cristo - Alexandre Dumas, đã sử dụng một hệ thống mã hóa màu sắc.

Có thể hiểu một cách đơn giản là trong mắt Alexandre mọi sự việc đều được ông quy về một màu sắc nào đó.

Các chuyên gia IT thường quy các điểm ảnh vi tính về hình ảnh thực tế: con chó, chiếc ô tô, cái cây,$\ldots$ thì ngược lại, nhà văn của chúng ta lại suy luận những sự vật thực tế về một màu sắc cụ thể. 

Ông cũng là người có nguyên tắc trong việc lựa chọn màu sắc cho tác phẩm của mình.

Quá thú vị đúng không?

Trong nhiều thập kỷ, thiên tài ấy đã sử dụng những màu sắc khác nhau cho những thể loại khác nhau của mình.

Màu xanh là màu của tiểu thuyết viễn tưởng.

Màu hồng dành cho những bài luận văn hay tác phẩm khoa học, còn màu vàng là để cho thơ.

Vậy tại sao lại không thử cách mã hóa này cho tác phẩm của bạn?

Hãy mạnh dạn thử qua mọi thứ trong cuộc sống này nhé!

\section{Tư thế dốc ngược}
Ấy là một phương pháp giúp chữa chứng ``đơ'' khi sáng tác -  ít nhất là tác giả nổi tiếng Dan Brown tin vào điều đó.

Chính ông đã thừa nhận rằng nó giúp ông thư giãn và tập trung hơn.

Treo ngược càng lâu, tác dụng càng lớn.
 
\section{Lặng nhìn một bức tường trống không}
Francine Prose, tác giả của cuốn tiểu thuyết Blue Angel tin rằng một bức tường trống là một phép ẩn dụ phù hợp cho tất cả những tác phẩm nghệ thuật.

Trong thời gian làm việc, Francine Prose đã đẩy hết những chiếc bàn ra phía cửa sổ, để trước mặt là một bức tường gạch.

Sự ``đơn điệu''  này giúp chúng ta khắc phục tình trạng mất tập trung và nhờ đó có thể làm việc hiệu quả trong nhiều giờ liền.

\section{Tự đối thoại}
Đã từng giành được nhiều giải thưởng, nhà biên kịch ``Mạng xã hội'' Aaron Sorkin thừa nhận từng bị gãy cả mũi trong quá trình làm việc.

Thật bất ngờ, chuyện đó thực sự đã xảy ra như thế nào?

Ông ấy thích các cuộc đối thoại của mình trước gương.

Có một lần vì quá nhập tâm mà ông đã tự đập đầu mình vào gương.

Kịch tính hóa là tốt, nhưng đừng quên sự an toàn.

\section{Lột bỏ bớt quần áo hoặc không mặc gì cả khi viết}
Bạn đang phải đối mặt với những deadline dày đặc, gặp nhiều bế tắc trong sáng tác.

Những khó khăn ấy chắc chắn sẽ bị đánh bại nếu bạn chú ý đến một phương pháp gây tò mò của nhà văn Victor Hugo.

Ông chỉ viết khi không mặc quần áo.

Có một khoảng thời gian Victor Hugo ngập ngụa trong lịch trình làm việc dày đặc để hoàn thiện cuốn tiểu thuyết kinh điển ``Nhà thờ Đức Bà Paris'', ông đã yêu cầu người hầu của mình giấu hết quần áo đi.

Điều này khiến ông không thể ra khỏi nhà.

Thậm chí là vào những ngày lạnh lẽo, Hugo cũng chỉ cho phép mình quấn một chiếc khăn trong lúc sáng tác.

Một hành động có chút điên rồ, quái dị.

\textsf{Nhà văn Victor Hugo.}

\section{Sự kích thích bằng cà-phê}
Tiểu thuyết gia người Pháp Honoré de Balzac đã nuôi dưỡng nguồn cảm hứng của mình bằng khoảng 50 tách cà-phê mỗi ngày.

Đây là lượng cà-phê đã truyền cảm hứng cho ông ở mức phù hợp.

Theo một số nghiên cứu, Balzac hầu như không ngủ trong khi viết cuốn The Human Comedy. Ngoài ra còn có một người cũng yêu thích cà-phê, đó là Voltaire.

Ông ấy cần tới 40 cốc mỗi ngày.\hfill$\square$

\begin{flushright}
    Nghiêm Anh | \url{Animedia.ru}
\end{flushright}

%------------------------------------------------------------------------------%

\chapter{\href{http://tramdoc.vn/tin-tuc/6-loi-ich-cua-thoi-quen-viet-moi-ngay-nzG1AW.html}{6 Lợi Ích của Thói Quen Viết Mỗi Ngày}}

\textbf{Tất cả chúng ta đều có thể giãi bày những suy nghĩ của mình lên mặt giấy. Điều đó có nghĩa rằng chúng ta đều là nhà văn, ngay cả khi không biết cách xoay chuyển cốt truyện một cách thuần thục như Tolstoy.}

%
Thói quen viết là một công cụ hữu ích để thể hiện bản thân, giúp phát triển khả năng sáng tạo và tư duy.

Và để có được điều đó, bạn hoàn toàn không cần phải ép buộc mình trở thành một tiểu thuyết gia bị ám ảnh với việc nhốt mình trong 4 bức tường.

Đơn giản chỉ cần biến việc viết trở thành những thói quen thân thuộc và từ từ, viết sẽ thay đổi cuộc sống của bạn.

\section{Viết là liệu trình tuyệt vời cho sức khỏe tinh thần}
Đa phần các nghiên cứu cho rằng, Expressive Writing (tạm dịch Văn biểu cảm - việc viết ra những gì bạn nghĩ và cảm nhận) sẽ mang lại cảm giác hạnh phúc cho người viết.

Một ví dụ rõ ràng nhất về Expressive Writing là viết nhật ký.

Điều này cũng được nhắc tới trong việc viết blog, nó giúp cải thiện sức khỏe tượng tự như viết tay vậy.

%
Ngoài ra, việc không thể diễn đạt suy nghĩ của mình bằng văn bản sẽ cản trở việc giao tiếp với người khác, nhất là trong việc trao đổi kinh nghiệm và cảm xúc.

Không dễ để cụ thể hóa những suy nghĩ của bạn và đưa chúng vào  một cuộc hội thoại.

Thói quen viết thường xuyên giúp đối phó với điều này.

%
Nghiên cứu của Laura King thuộc Đại học Southern Methodist cho thấy những người viết về mục tiêu, ước mơ và thành tích của họ thường hạnh phúc và khỏe mạnh hơn.

Tôi cũng nhận thấy tác dụng tương tự: những người bị căng thẳng trong công việc bắt đầu thói quen ghi chép lại trong vài ngày, sau đó cảm thấy tốt hơn và năng suất của họ tăng 29\%.

(Adam Grant, nhà văn, nhà báo, giáo sư tại Trường Kinh doanh Wharton).

%
Bên cạnh đó, thói quen viết lách giúp bạn dễ dàng truyền đạt ngay cả những ý tưởng phức tạp cho người khác.

Nó giúp loại bỏ lời biện hộ ``Ở trong đầu có vẻ tốt hơn''.

Dĩ nhiên là, khi viết bạn phải trình bày rõ ràng suy nghĩ của mình.

\section{Thói quen viết giúp bạn vượt qua những giai đoạn khó khăn}
Trong một nghiên cứu được thực hiện giữa các kỹ sư bị sa thải gần đây, người ta thấy rằng những người thường xuyên viết và truyền đạt tâm tư của mình trên giấy sẽ nhanh chóng tìm được việc làm mới.

%
Các kỹ sư ghi lại suy nghĩ và cảm xúc của họ về việc bị sa thải ít cảm thấy tức giận và thù địch hơn đối với người chủ cũ.

Họ cũng uống ít hơn.

Kết quả là, 8 tháng sau, 52\% kỹ sư, những người thường xuyên viết, đã tìm được công việc toàn thời gian mới, so với chỉ 19\% ở nhóm được yêu cầu ``không viết''.

(Theo Adam Grant, nhà văn, nhà báo, giáo sư tại Trường Kinh doanh Wharton).

%
Những người tham gia thử nghiệm nói rằng: khi mô tả những tâm tư mà mình  không thể chia sẻ với ai, họ không phủ nhận những khó khăn, mà cố gắng chấp nhận và vượt qua chúng.

Do đó, theo thời gian, những lo lắng dần phai nhạt.

%
Việc viết về những sự kiện đau buồn giúp bạn giải phóng cảm giác chán nản.

Tuy nhiên, bạn sẽ mất ít nhất 6 tháng để trải nghiệm toàn bộ lợi ích của hoạt động này.

Đồng thời, để thoát khỏi những lo lắng tiêu cực, bạn không nên ép buộc bản thân.

Việc viết lách, ghi chép phải tự nhiên, phải mang đến sự hài lòng cho người thực hiện công việc đó.

\section{Thói quen viết giúp bạn lấy lại động lực}
Ở một nghiên cứu khác cho rằng, đối với những đối tượng viết ra ít nhất một lần một tuần, họ nhận ra có nhiều thứ tốt đẹp đang diễn ra trong cuộc sống và họ nhìn vào tương lai với phong thái lạc quan, có động lực nhiều hơn.

%
Nhưng có một ``ngoại lệ'': nếu bạn viết mỗi ngày, thì sẽ không có sự khác biệt đáng kể.

Điều này có nghĩa rằng: bất kỳ việc gì, nếu bạn làm việc đó quá thường xuyên mà chẳng có một mục đích thực sự, mong muốn thực sự, thì sự chán nản sẽ nhanh chóng kéo đến thôi.

\section{Khi viết, hãy cố gắng sắp xếp mọi thứ trong đầu}
Bạn đã bao giờ mở nhiều tab (cửa sổ trình duyệt web) trong trình duyệt của mình cùng một lúc chưa?

Thật khó để không bị nhầm lẫn và không bị chúng làm phân tâm.

Ở trong đầu cũng vậy, cố gắng suy nghĩ nhiều ý tưởng, hành động, kế hoạch cùng một lúc cũng chẳng khác nào mở những tab giống nhau.

Thói quen viết sẽ hình thành suy nghĩ của bạn, bạn chuyển chúng từ não ra giấy và giải phóng tâm trí mình.

\textsf{Thói quen viết sẽ hình thành suy nghĩ của bạn.}

\section{Viết giúp ích cho trí nhớ}
Thông tin sẽ dễ nhớ hơn khi bạn hiểu về tầm quan trọng của chúng và diễn đạt chúng bằng ngôn từ của mình.

Có được 1 bài viết thú vị đòi hỏi bạn phải có nguyên tắc và sự tự chủ trong tổ chức: cần phải thường xuyên tập trung, tìm kiếm các nguồn thông tin mới, cảm hứng và kiến thức mới.

%
Khi tìm kiếm những ý tưởng mới, bạn sẽ phát triển được tư duy, khả năng phân tích và nghiên cứu của mình, học cách đi sâu và tìm ra những chủ đề mà bạn quan tâm.

Bằng cách dành thời gian để viết, bạn học được cách giải quyết vấn đề hiệu quả hơn.

%
Trong 1 khoảng thời gian, nếu viết cùng một chủ đề, bạn sẽ nhanh chóng chuyển từ những ý tưởng đại trà sang những ý tưởng mới mẻ, và sau đó bạn có thể tạo ra gì đó độc đáo.

Vì vậy, nhiều nhà văn bắt đầu với một đoạn văn, sau đó chuyển thành một bài tiểu luận, sau đó là một bài luận rồi đến một loạt các bài báo, và các bài báo trở thành một cuốn sách.

\section{Thói quen viết giúp bạn chấp nhận những lời chỉ trích}
Trong thế giới hiện đại, mọi người đều muốn làm cho mình được biết đến bằng cách này hay cách khác.

Mỗi người có thể xuất bản tác phẩm của mình và chia sẻ nó với những người khác.

%
Hãy nghĩ mà xem, thật tuyệt vời khi bạn có thể ảnh hưởng đến người khác bằng ngôn từ của mình.

Khi ai đó viết cho bạn một lá thư cảm ơn về điều mà bạn đã chia sẻ, bạn có thể sẽ rất ngạc nhiên đấy.

%
Và khi đối mặt với những lời chỉ trích, các nhà văn sẽ trở nên kiên cường hơn.

Những lời nhận xét, ngay cả khi chúng phi lý thì đó cũng là một cách tôi luyện tuyệt vời.

\begin{flushright}
    Nghiêm Anh | Theo \url{https://animedia-company.cz/writing-regularly-6-benefits/}
\end{flushright}

%------------------------------------------------------------------------------%



%------------------------------------------------------------------------------%

\printbibliography[heading=bibintoc]
\end{document}